\section{Related Work}
\label{sec:related}

Unsupervised time-series anomaly detection spans dynamical equation discovery, reconstructive generative modeling, contrastive representation learning, and predictive world models. This section reviews these approaches and positions DyFlow-JEPA.

\subsection{Continuous Dynamical Systems and Flow Matching}
The physical modeling of multi-sensor telemetry has long sought to capture system behavior as continuous dynamical trajectories governed by differential equations $\dot{x}(t) = \mathbf{f}(x(t))$ \cite{brunton2016discovering}.
\begin{itemize}[leftmargin=*]
    \item \textit{Neural Ordinary Differential Equations (Neural ODEs) \cite{chen2018neural}:} Parameterize continuous velocity fields using neural networks $\frac{dx}{dt} = f_\theta(x(t), t)$ and solve trajectories via numerical ODE integrators. While Neural ODEs learn flexible vector fields, fitting them directly in raw sensor space forces numerical solvers to integrate through high-frequency stochastic sensor noise, leading to numerical stiffness, slow training, and extreme sensitivity to measurement noise.
    \item \textit{Conditional Flow Matching (CFM) \cite{lipman2022flow, tong2023improving}:} Continuous Normalizing Flows and Optimal Transport Conditional Flow Matching (OT-CFM) train neural velocity vector fields $v_\psi(z_t, t)$ along straight probability paths. Stochastic interpolant frameworks \cite{albergo2022building} provide efficient simulation-free training objectives. DyFlow-JEPA couples Optimal Transport Flow Matching with non-contrastive Joint-Embedding representations, learning continuous dynamical vector fields purely in regularized latent space without raw-space noise integration.
\end{itemize}

\subsection{Reconstruction-Based Detection}
The predominant deep learning approach in time-series anomaly detection relies on learning nominal patterns through raw signal reconstruction. Hundman et al. \cite{Hundman2018b} introduced LSTM-based prediction coupled with non-parametric dynamic thresholding for NASA spacecraft telemetry (SMAP/MSL). Modern architectures incorporate specialized attention and frequency-domain transformations:
\begin{itemize}[leftmargin=*]
    \item \textit{Anomaly Transformer \cite{xu2022anomaly}:} Introduces Anomaly-Attention to quantify the \textit{Association Discrepancy} between local prior Gaussian kernel associations and global series self-attention. While effective for isolated point anomalies, it couples discrepancy with point-level waveform reconstruction.
    \item \textit{TimesNet \cite{wu2023timesnet}:} Decomposes 1D multivariate time series into 2D temporal variations using Fast Fourier Transform (FFT) periodicity discovery and 2D Inception blocks. However, optimizing point-level 2D reconstruction makes it susceptible to over-fitting high-frequency sensor noise.
    \item \textit{TranAD \cite{tuli2022tranad}:} Employs a two-phase adversarial Transformer with sub-series attention and focus scoring to amplify subtle reconstruction errors, but remains vulnerable to phase-drift in continuous dynamical systems.
\end{itemize}

Reconstructive methods share a common vulnerability: \textit{signal-level reconstruction memorization}. Because loss functions optimize point-level $L_1$ or $L_2$ differences in raw signal space, reconstructive networks are biased toward fitting high-frequency stochastic noise rather than learning invariant dynamical laws. Furthermore, during sustained anomalies, corrupted observations contaminate the context window, causing the reconstructor to adapt and reconstruct the anomaly as the "new normal."

\subsection{Contrastive Learning and Representation-Based Detection}
To avoid reconstructing raw noise, contrastive methods learn discriminative representations by separating nominal sequences from anomalous perturbations:
\begin{itemize}[leftmargin=*]
    \item \textit{Neural Contextual Anomaly Detection (NCAD) \cite{carmona2021neural}:} Uses Temporal Convolutional Networks (TCNs) trained with Contextual Outlier Exposure (COE). Artificial synthetic anomalies (Gaussian noise, point spikes, trend shifts) are injected into training windows to create negative pairs. However, synthetic injections distort the physical dynamical manifold and fail to generalize to subtle real-world physical faults (e.g., gradual thermal dissipation or gait freezing).
    \item \textit{DCdetector \cite{yang2023dcdetector}:} Replaces reconstruction with multi-scale dual-attention (patch-wise and channel-wise) contrastive representation learning, maximizing cross-scale representation consistency. While DCdetector eliminates artificial injection heuristics, contrastive objectives require careful scale selection and representation pooling.
\end{itemize}

\subsection{Non-Contrastive Learning and Time-Series JEPAs}
Joint-Embedding Predictive Architectures (JEPA), originally proposed by LeCun \cite{lecun2022path} and demonstrated in computer vision by I-JEPA \cite{assran2023self} and V-JEPA \cite{bardes2024vjepa}, predict target representations from context in latent space. Non-contrastive regularizers like Barlow Twins \cite{zbontar2021barlow}, W-MSE \cite{ermolov2021whitening}, and VICReg \cite{bardes2021vicreg} prevent representation collapse without negative mining.

Recent works have extended joint-embedding predictive principles to time-series sequences:
\begin{itemize}[leftmargin=*]
    \item \textit{Time-Series JEPAs:} SC-JEPA / MTS-JEPA \cite{he2026scjepa} stabilizes latent predictive learning for time-series anomaly prediction using a multi-scale soft-codebook dictionary. T-SAR-JEPA \cite{woldesenbet2026tsarjepa} applies self-supervised latent prediction to Synthetic Aperture Radar (SAR) amplitude time series. CF-JEPA \cite{lee2026cfjepa} employs mask-free forward prediction with asymmetric encoders for time-series representation learning. In cyber-physical and automotive domains, Girgis et al. \cite{girgis2026timeseries} explore time-series JEPAs for predictive remote control over capacity-constrained IoT channels, while Fertig et al. \cite{fertig2026online} investigate JEPA embeddings for online automotive sensor monitoring.
\end{itemize}
However, these time-series JEPA adaptations rely on deterministic point regression or discrete codebook quantization, which collapse to the conditional mean on stochastic, multi-branching physical systems.

\subsection{Flow Matching in Joint Embeddings (Vision, Control, and World Models)}
Recently, the integration of continuous flow matching into joint-embedding architectures has been applied to generative vision and reinforcement learning:
\begin{itemize}[leftmargin=*]
    \item \textit{Generative Vision and Multimodal Synthesis:} D-JEPA \cite{xie2024djepa} pioneered the use of flow matching and diffusion objectives within JEPA predictors to replace discrete next-token prediction for continuous image, video, and audio generation. Concurrently, JEPA-T \cite{wang2025jepat} integrated flow-matching feature predictors with cross-attention for text-to-image synthesis.
    \item \textit{World Models and Robotic Control:} In reinforcement learning, Qantara \cite{qantara2026} introduced Bridge-Flow training for multi-task JEPA control, while concurrent work by Madaan et al. \cite{flowjepa2026world} investigated flow matching for trajectory-level state rollouts in latent dynamics world models.
\end{itemize}

Prior work applies flow matching to generative image synthesis and action-conditioned robotic state transitions. In contrast, DyFlow-JEPA applies continuous flow matching within joint embeddings to unsupervised multivariate time-series anomaly detection. Rather than optimizing action-conditioned rollouts or pixel generation, DyFlow-JEPA models the continuous tangent velocity field of physical dynamical attractors in multi-sensor telemetry, evaluating it with a deterministic midpoint discrepancy estimator, Mahalanobis mode whitening, and Extreme Value Theory (EVT / SPOT) tail calibration.

\subsection{Positioning and Method Design}
DyFlow-JEPA combines continuous dynamical modeling and joint-embedding self-supervised learning through four key components:
\begin{enumerate}[leftmargin=*]
    \item \textit{Continuous Latent Vector Field Modeling:} Instead of raw-space differential equation solving or deterministic point regression, DyFlow-JEPA learns the continuous tangent velocity field of the physical manifold via Optimal Transport Conditional Flow Matching directly in regularized latent space.
    \item \textit{Non-Contrastive Manifold Regularization:} Replaces synthetic negative injection heuristics with non-contrastive variance and spectral covariance penalties, precluding dimensional collapse while preserving smooth attractor geometry.
    \item \textit{Deterministic Midpoint Velocity Discrepancy:} Evaluates the velocity field at the Optimal Transport midpoint ($t=0.5$), providing an efficient, zero-Monte-Carlo-variance proxy for trajectory deviation whitened by an empirical Mahalanobis precision tensor.
    \item \textit{Extreme Value Tail Calibration:} Applies Extreme Value Theory (EVT / SPOT) \cite{siffer2017anomaly} under the Pickands-Balkema-de Haan theorem \cite{balkema1974residual} to fit Generalized Pareto tails on nominal velocity residuals, setting decision boundaries without manual threshold tuning.
\end{enumerate}

